\documentclass[11pt]{article}
\usepackage[margin=1.05in]{geometry}
\usepackage{amsmath,amssymb,amsthm,mathtools}
\usepackage{enumitem}
\usepackage{microtype}
\usepackage[colorlinks=true,linkcolor=blue,citecolor=blue,urlcolor=blue]{hyperref}

\newtheorem{theorem}{Theorem}[section]
\newtheorem{proposition}[theorem]{Proposition}
\newtheorem{lemma}[theorem]{Lemma}
\newtheorem{corollary}[theorem]{Corollary}
\newtheorem{inputtheorem}[theorem]{Input theorem}
\theoremstyle{remark}
\newtheorem{remark}[theorem]{Remark}

\newcommand{\Pp}{\mathbb P}
\newcommand{\Ee}{\mathbb E}
\newcommand{\one}{\mathbf 1}
\newcommand{\dd}{\,d}
\newcommand{\Exp}{\operatorname{Exp}}
\newcommand{\ord}{\operatorname{ord}}
\newcommand{\lcm}{\operatorname{lcm}}
\newcommand{\Cyc}{\operatorname{Cyc}}
\newcommand{\Var}{\operatorname{Var}}

\title{An Erd\H{o}s--Tur\'an Law for Comparable-Rate Luce Permutations}
\author{Christopher D. Long}
\date{}

\begin{document}
\maketitle

\begin{abstract}
Let $\sigma_n$ be the Luce permutation induced by independent exponential
clocks whose rates, after a common rescaling, satisfy
$1\leq\theta_{n,i}\leq\kappa$ for one fixed $\kappa<\infty$.  We prove
\[
 \frac{\log\ord(\sigma_n)-\frac12(\log n)^2}
 {\sqrt{\frac13(\log n)^3}}
 \Rightarrow N(0,1).
\]
Thus the classical Erd\H{o}s--Tur\'an law persists for arbitrary
comparable-rate triangular arrays; no limiting rate profile is required.
The proof treats the additive proxy $\sum_r C_{n,r}\log r$ by a weighted
clock-swap characteristic-function argument.  A quantitative
finite-endpoint comparison is used through a theorem-level interface.  The
arithmetic difference between the product and the least common multiple of
the cycle lengths is controlled by an exact rate-oriented three-clock
switch, whose orientation leaves only one positive coefficient in the
Radon--Nikodym logarithm and yields a uniform two-cycle factorial-moment
bound.
\end{abstract}

\tableofcontents

\section{Introduction}

For a permutation $\pi\in S_n$, its order is the least positive integer
$k$ for which $\pi^k$ is the identity, equivalently
\[
 \ord(\pi)=\lcm\{\text{cycle lengths of }\pi\}.
\]
For a uniformly chosen permutation, Erd\H{o}s and Tur\'an proved that the
logarithm of the order is asymptotically Gaussian on the
$(\log n)^{3/2}$ scale~\cite{ErdosTuran1967}.  Sharp rates and a broader
probabilistic theory were subsequently developed for the uniform and Ewens
settings and for logarithmic combinatorial structures; see, among others,
\cite{BarbourTavare1994,ABT2003,GnedinIksanovMarynych2012,StormZeindler2015}.

The Luce distribution is generated by ranking independent exponential
clocks with prescribed rates.  It is the stationary distribution of the
Tsetlin library and is typically nonexchangeable and spatially
inhomogeneous~\cite{Luce1959,CDK2024}.  Recent work developed the permuton,
local, and multiscale cycle theories of this model
\cite{BCD2025,LongLuce2026,LongSwitch2026}.  The permutation order is a
particularly sensitive remaining statistic: it depends on divisibility
relations among cycle lengths across fixed, mesoscopic, and macroscopic
scales, and is not determined by a permuton limit or by macroscopic
Poisson--Dirichlet convergence alone.

Our result shows that the classical order law is nevertheless universal
throughout the comparable-rate class.

\begin{theorem}[Main theorem]\label{thm:main}
Let $\sigma_n$ be a Luce permutation with positive rates
$\theta_{n,1},\ldots,\theta_{n,n}$ satisfying
\[
 \sup_n\frac{\max_i\theta_{n,i}}{\min_i\theta_{n,i}}<\infty.
\]
Then
\[
 \frac{\log\ord(\sigma_n)-\frac12(\log n)^2}
 {\sqrt{\frac13(\log n)^3}}
 \Rightarrow N(0,1).
\]
\end{theorem}

The theorem is stronger in hypothesis scope than a sampled-profile result:
the rates may form an arbitrary triangular array, and neither convergence
nor regularity of a macroscopic profile is assumed.

\subsection*{Proof architecture}

After a harmless common rescaling, take $1\leq\theta_{n,i}\leq\kappa$.
The argument has four parts.

\begin{enumerate}[label=\textnormal{(\roman*)}]
\item The product of the cycle lengths has logarithm
\[
 B_n=\sum_{r=1}^n C_{n,r}\log r.
\]
A weighted clock-swap identity and a quantitative finite-endpoint
comparison yield the same characteristic exponent as an independent
Poisson cycle process, up to an error negligible on the
$(\log n)^{3/2}$ scale.

\item The quantitative finite-endpoint theory is imported only through the
precise Input Theorem~\ref{input:endpoint}.  All subsequent summations and
optimizations are carried out in this paper.

\item A new rate-oriented three-clock change of measure controls two cycle
counts simultaneously.  The cyclic orientation is chosen by decreasing
rate, which makes the Radon--Nikodym logarithm have exactly one positive
clock coefficient and gives a uniform $C_\kappa/(rs)$ factorial-moment
bound.

\item The von Mangoldt decomposition converts the difference between the
cycle product and the cycle least common multiple into prime-power overlap
counts.  The one- and two-cycle estimates imply an $O_\kappa(\log n\log\log
n)$ bound in mean, which is negligible at the Erd\H{o}s--Tur\'an scale.
\end{enumerate}

\subsection*{Dependencies and quantitative refinements}

The only imported probabilistic inputs are the marked one-switch estimates
in Input Theorem~\ref{input:one-switch} and the quantitative
finite-endpoint interface in Input Theorem~\ref{input:endpoint}, both from
the companion switching manuscript~\cite{LongSwitch2026}.  The weighted
characteristic-function argument, the rate-oriented three-clock identity,
and the arithmetic product--LCM comparison are proved here.  A separate
quantitative supplement~\cite{LongBerry2026} proves the sharp
$O_\kappa((\log n)^{-1/2})$ Kolmogorov rate under exact-mean centering and a
first Edgeworth expansion.  That argument uses a stronger linear-depth form
of the finite-endpoint comparison and fixed-order multi-clock estimates;
neither strengthening is needed for the qualitative theorem proved here.

\section{Model, switching inputs, and cutoffs}

Let
\[
 X_{n,i}\sim\Exp(\theta_{n,i}),\qquad 1\leq i\leq n,
\]
be independent, with
\begin{equation}\label{eq:comparable}
 1\leq\theta_{n,i}\leq\kappa.
\end{equation}
Let $\sigma_n(i)$ be the rank of $X_{n,i}$ among the $n$ clocks.  Write
$C_{n,r}$ for the number of $r$-cycles of $\sigma_n$ and $L_n(i)$ for the
length of the cycle containing $i$.

For distinct labels $i,j$, put
\[
 R_{ij}=\exp\{-(\theta_i-\theta_j)(X_j-X_i)\},
 \qquad R_{ii}=1,
\]
and define the switching intensity
\[
 S_{n,r}=\sum_{i:L_n(i)>r}R_{i,\sigma_n^r(i)}.
\]
We import the following estimates from the clock-swap and finite-endpoint
mixing theory in~\cite{LongSwitch2026}.

\begin{inputtheorem}[Marked one-switch identity, moments, and tails]\label{input:one-switch}
For every $1\leq r<n$ and every bounded complex-valued function
$\Phi:S_n\to\mathbb C$,
\begin{equation}\label{eq:marked-one-switch}
 r(n-r)\Ee\bigl[C_{n,r}\Phi(\sigma_n)\bigr]
 =\Ee\sum_{i:L_n(i)>r}
 R_{i,\sigma_n^r(i)}
 \Phi\bigl(\sigma_n\circ(i\,\sigma_n^r(i))\bigr).
\end{equation}
Taking $\Phi\equiv1$ gives
\begin{equation}\label{eq:one-switch}
 r(n-r)\Ee C_{n,r}=\Ee S_{n,r}.
\end{equation}
Set
\[
 \lambda_\kappa=
 \begin{cases}
  \infty,&\kappa=1,\\[1mm]
  \dfrac{\kappa}{\kappa-1},&\kappa>1.
 \end{cases}
\]
For every $1<p<\lambda_\kappa$,
\begin{equation}\label{eq:selected-p}
 \sup_{n,r}\frac1n\Ee\sum_{i:L_n(i)>r}
 R_{i,\sigma_n^r(i)}^p<\infty.
\end{equation}
If $\kappa>1$, then, uniformly for $A\geq1$,
\begin{align}
 \sup_{n,r}\frac1n\Ee\sum_{i:L_n(i)>r}
 R_{i,\sigma_n^r(i)}\one_{\{R_{i,\sigma_n^r(i)}>A\}}
 &\leq C_\kappa A^{-1/(\kappa-1)}\label{eq:selected-tail}\\
 \sup_{n}\max_{i,j}\Ee\left[
 R_{ij}\one_{\{R_{ij}>A\}}\right]
 &\leq C_\kappa A^{-1/(\kappa-1)}.\label{eq:fresh-tail}
\end{align}
Moreover, $\Ee R_{ij}=1$ for every $i,j$.  If $\kappa=1$, all the factors
$R_{ij}$ are identically one.
\end{inputtheorem}

In particular, $\Ee S_{n,r}\leq C_\kappa n$ uniformly in $n,r$.

\begin{lemma}[Crude cycle, root, and logarithmic-tail bounds]
\label{lem:crude}
Uniformly in $1\leq r<n$,
\begin{equation}\label{eq:Cr-crude}
 \Ee C_{n,r}\leq \frac{C_\kappa n}{r(n-r)}.
\end{equation}
If $I_n$ is uniform on $[n]$, independently of the clocks, then, for
$a\leq n/2$,
\begin{equation}\label{eq:root-crude}
 \Pp\{L_n(I_n)\leq a\}\leq C_\kappa a/n.
\end{equation}
Consequently,
\begin{equation}\label{eq:log-cycle-tail}
 \Ee\sum_{r\leq b}(\log r)C_{n,r}=O_\kappa((\log b)^2),
\end{equation}
and, for $U\leq n/2$,
\begin{equation}\label{eq:large-log-tail}
 \Ee\sum_{r>U}(\log r)C_{n,r}
 \leq C_\kappa\log n\{1+\log(n/U)\}.
\end{equation}
\end{lemma}

\begin{proof}
Equation~\eqref{eq:Cr-crude} follows from~\eqref{eq:one-switch} and
$\Ee S_{n,r}\leq C_\kappa n$.  Since
\[
 \Pp\{L_n(I_n)=r\}=\frac{r\Ee C_{n,r}}n,
\]
for $a\leq n/2$ we have
\[
 \Pp\{L_n(I_n)\leq a\}
 \leq C_\kappa\sum_{r\leq a}\frac1{n-r}
 \leq C_\kappa'\frac an.
\]
For $r\leq n/2$, equation~\eqref{eq:Cr-crude} gives
$\Ee C_{n,r}\leq 2C_\kappa/r$.  Summing $(\log r)/r$ proves
\eqref{eq:log-cycle-tail}.  For $U<r\leq n/2$, the same estimate gives
\[
 \Ee\sum_{U<r\leq n/2}(\log r)C_{n,r}
 \leq C_\kappa\log n\sum_{U<r\leq n/2}\frac1r.
\]
There is deterministically at most one cycle of length greater than $n/2$,
which contributes at most $\log n$.  This proves~\eqref{eq:large-log-tail}.
\end{proof}

Fix once and for all a number $p$ satisfying
\[
 1<p<\lambda_\kappa
\]
(with any $p>1$ when $\kappa=1$), and set
\[
 \beta=1-\frac1p>0.
\]
Choose a fixed constant
\begin{equation}\label{eq:A0-choice}
 A_0>\max\left\{2,\frac1{2\beta}\right\},
\end{equation}
and put
\begin{equation}\label{eq:cutoffs}
 b_n=\left\lceil(\log n)^3\right\rceil,
 \qquad
 U_n=\left\lfloor\frac{n}{(\log n)^{A_0}}\right\rfloor.
\end{equation}
For all sufficiently large $n$,
\begin{equation}\label{eq:cutoff-range}
 2U_n\leq n/2.
\end{equation}

\section{Quantitative switching on dyadic blocks}

Put
\[
 \Delta_{n,r}=\Ee\left|\frac{S_{n,r}}n-1\right|.
\]
Fix an integer $a$ with $b_n\leq a\leq U_n$.  In the two-root endpoint
argument take lower length $a$, upper length $2a$, and hence
\begin{equation}\label{eq:block-parameters}
 \delta=\frac{4a}{n},\qquad
 \ell=\log\frac1\delta,\qquad
 h=\delta^{1/4},\qquad
 \rho=\delta^{1/(4\kappa)},\qquad
 M=n\delta^{1/2},
\end{equation}
and
\begin{equation}\label{eq:m-def}
 m=\left\lfloor\min\left\{\frac a3,\sqrt\ell\right\}\right\rfloor.
\end{equation}
These are the parameters of the finite-endpoint construction with two
independent roots and maximal exposed length $2a$.

For $A\geq1$, define
\[
 F_A(i,j)=R_{ij}\wedge A,
 \qquad
 S_{n,r}^{(A)}=\sum_{i:L_n(i)>r}F_A(i,\sigma_n^r(i)),
\]
and
\[
 D_{n,A}=\frac1{n^2}\sum_{i,j=1}^nF_A(i,j),
 \qquad
 \mu_{n,A}=\Ee D_{n,A}.
\]

The following is the complete interface from the finite-endpoint theory
used in this paper.  It is the quantitative form of the one- and two-root
endpoint comparison proved in Sections~6--8 of~\cite{LongSwitch2026}.

\begin{inputtheorem}[Quantitative finite-endpoint comparison]
\label{input:endpoint}
There are constants $c,C,\alpha_0>0$, depending only on $\kappa$, such
that, uniformly for $b_n\leq a\leq U_n$, $a\leq r<2a$, and $A\geq1$,
\begin{align}
 \left|\Ee\frac{S_{n,r}^{(A)}}n-\mu_{n,A}\right|
 &\leq C_\kappa A\Gamma_{n,a},\label{eq:first-A}\\
 \left|\Ee\left(\frac{S_{n,r}^{(A)}}n\right)^2
       -\Ee D_{n,A}^2\right|
 &\leq C_\kappa A^2\Gamma_{n,a},\label{eq:second-A}
\end{align}
where
\begin{align}
 \Gamma_{n,a}\leq C\bigg[&
 (1-\alpha_0)^{m-1}
 +m\delta^{1/4}
 +\frac{m^2}{n\delta^{1/2}}
 +m\delta^{1/(4\kappa)}
 +\frac{m\log n}{n\delta^{1/2}}
 \notag\\
 &+\delta^{1/(4\kappa)}
 +\delta^{1/2}\ell
 +\delta
 +g_{n,a}\bigg],
 \label{eq:Gamma}
\end{align}
with
\begin{equation}\label{eq:g-explicit}
 g_{n,a}\leq C\left[
 \delta^{-1/4}\ell\,e^{-c n\delta^{1/2}}
 +n^2\delta^{-1/4}e^{-c n\delta^{3/4}}
 +n^{-2}
 \right].
\end{equation}
The constants are independent of the truncation level $A$.
\end{inputtheorem}

\begin{remark}[Use of the input theorem]
Only the bounds \eqref{eq:first-A}--\eqref{eq:g-explicit} are used below.
In particular, the clock-block construction, conditional permanental
matching, reset chain, and adaptive exposure bookkeeping from the companion
paper do not enter the present proof except through this stated interface.
The companion endpoint theorem admits a free terminal-depth parameter.  On
the present dyadic range, this paper uses the square-root-depth specialization
$m\asymp\sqrt\ell$, which suffices for the qualitative
characteristic-function argument, whereas the quantitative
supplement~\cite{LongBerry2026} uses $m\asymp\ell$ to obtain polynomial
product--Poisson accuracy.
\end{remark}

\begin{proposition}[One-block quantitative switching]\label{prop:block}
There is $C_\kappa<\infty$ such that, uniformly for
$b_n\leq a\leq U_n$ and $a\leq r<2a$,
\begin{equation}\label{eq:block-delta}
 \Delta_{n,r}\leq C_\kappa(\Gamma_{n,a}+n^{-1})^{1/(2\kappa)}.
\end{equation}
\end{proposition}

\begin{proof}
Replacing one clock $X_k$ changes only the terms of $D_{n,A}$ with
$i=k$ or $j=k$, at most $2n$ terms, each lying in $[0,A]$.  Hence
$|D_{n,A}-D_{n,A}^{(k)}|\leq2A/n$, and the Efron--Stein inequality gives
\begin{equation}\label{eq:D-var}
 \Var(D_{n,A})\leq\frac{4A^2}{n}.
\end{equation}
Put $X_{n,r}^{(A)}=S_{n,r}^{(A)}/n$.  Since
$0\leq\mu_{n,A}\leq1$,
\begin{align*}
 \Ee|X_{n,r}^{(A)}-\mu_{n,A}|^2
 ={}&\bigl(\Ee(X_{n,r}^{(A)})^2-\Ee D_{n,A}^2\bigr)
 +\Var(D_{n,A})\\
 &-2\mu_{n,A}\bigl(\Ee X_{n,r}^{(A)}-\mu_{n,A}\bigr).
\end{align*}
Equations~\eqref{eq:first-A}, \eqref{eq:second-A}, and~\eqref{eq:D-var}
therefore imply
\begin{equation}\label{eq:truncated-L2-quant}
 \Ee|X_{n,r}^{(A)}-\mu_{n,A}|^2
 \leq C_\kappa A^2(\Gamma_{n,a}+n^{-1}).
\end{equation}

Suppose first that $\kappa>1$.  Equations~\eqref{eq:selected-tail} and
\eqref{eq:fresh-tail}, together with $\Ee R_{ij}=1$, give
\[
 \Ee\frac{S_{n,r}-S_{n,r}^{(A)}}n
 +|1-\mu_{n,A}|
 \leq C_\kappa A^{-1/(\kappa-1)}.
\]
By~\eqref{eq:truncated-L2-quant},
\begin{equation}\label{eq:Delta-opt}
 \Delta_{n,r}
 \leq C_\kappa\left[
 A(\Gamma_{n,a}+n^{-1})^{1/2}
 +A^{-1/(\kappa-1)}
 \right].
\end{equation}
For $\Gamma_{n,a}+n^{-1}\leq1$, choose
\[
 A=(\Gamma_{n,a}+n^{-1})^{-(\kappa-1)/(2\kappa)};
\]
when the quantity exceeds one, the claimed bound follows after enlarging
the constant.  This proves~\eqref{eq:block-delta}.  If $\kappa=1$, all swap
factors equal one, no truncation is needed, and
\eqref{eq:truncated-L2-quant} with $A=1$ gives the same conclusion with
exponent $1/2=1/(2\kappa)$.
\end{proof}

\begin{corollary}[Logarithmically weighted dyadic summability]
\label{cor:weighted-sum}
Let $a_j=2^jb_n$.  Then
\begin{equation}\label{eq:weighted-sum}
 \sum_{j:\,a_j\leq U_n}
 \log(2a_j)
 \sup_{a_j\leq r<2a_j}\Delta_{n,r}
 =o((\log n)^{3/2}).
\end{equation}
\end{corollary}

\begin{proof}
Set
\[
 \eta=\frac1{2\kappa},
 \qquad
 \delta_j=\frac{4a_j}{n},
 \qquad
 L_j=\log\frac1{\delta_j}.
\]
Since $a_j\geq b_n=(\log n)^3$ and $L_j\leq\log n$, the minimum in
\eqref{eq:m-def} is attained by $\sqrt{L_j}$ for all sufficiently large
$n$, uniformly in $j$.  Thus $m_j\asymp\sqrt{L_j}$.  Moreover,
\[
 L_j\geq A_0\log\log n-O(1).
\]
Since $0<\eta\leq1$, the map $t\mapsto t^\eta$ is subadditive on
$[0,\infty)$.  We may therefore sum the terms in~\eqref{eq:Gamma}
separately after applying Proposition~\ref{prop:block}.

For the reset-mixing term,
\begin{align*}
 \sum_j\log(2a_j)e^{-c\eta m_j}
 &\leq C\log n\sum_{L_j\geq A_0\log\log n-O(1)}e^{-c'\sqrt{L_j}}\\
 &\leq C\log n\,(1+\sqrt{\log\log n})
 e^{-c''\sqrt{\log\log n}}
 =o((\log n)^{3/2}).
\end{align*}

After taking the $\eta$-power, each term involving a positive power of
$\delta_j$ is bounded by
\[
 C L_j^d e^{-\gamma L_j}
\]
for fixed $d\geq0$ and $\gamma>0$.  Since the $L_j$ form an arithmetic
lattice with fixed spacing,
\[
 \sum_j L_j^d e^{-\gamma L_j}
 \leq C(\log\log n)^d(\log n)^{-\gamma A_0}.
\]
Multiplication by the weight $\log(2a_j)\leq\log n$ is still
$o((\log n)^{3/2})$.  This treats
$m\delta^{1/4}$, $m\delta^{1/(4\kappa)}$,
$\delta^{1/(4\kappa)}$, $\delta^{1/2}\ell$, and $\delta$.

For the terms containing $M^{-1}=(n\delta_j^{1/2})^{-1}$, use
$m_j\leq\sqrt{\log n}$ and $a_j\geq(\log n)^3$ to obtain
\[
 \frac{m_j^2+m_j\log n}{n\delta_j^{1/2}}
 \leq C\frac{(\log n)^{3/2}}{\sqrt{na_j}}.
\]
Since $a_j\geq(\log n)^3$, the apparent factor
$(\log n)^{3/2}$ is cancelled by $\sqrt{a_j}$.  After taking the
$\eta$-power, multiplying by a logarithmic weight, and summing over
$O(\log n)$ blocks, the contribution is therefore at most
\[
 C(\log n)^2
 \left(\frac1{\sqrt n}\right)^\eta=o(1).
\]

The first two terms of~\eqref{eq:g-explicit} are superpolynomially small
uniformly for $a_j\geq(\log n)^3$.  The $n^{-2}$ term and the additional
$n^{-1}$ in~\eqref{eq:block-delta} contribute at most
$C(\log n)^2n^{-\eta}=o(1)$.  Combining the estimates proves\eqref{eq:weighted-sum}.
\end{proof}

\section{A central limit theorem for the product of the cycle lengths}

Define
\begin{equation}\label{eq:B-def}
 B_n=\sum_{r=1}^n(\log r)C_{n,r}
 =\log\prod_{C\in\Cyc(\sigma_n)}|C|,
\end{equation}
and its truncated version
\begin{equation}\label{eq:Bcirc-def}
 B_n^\circ=\sum_{b_n<r\leq U_n}(\log r)C_{n,r}.
\end{equation}

\begin{lemma}[Truncation of the additive proxy]\label{lem:B-trunc}
One has
\begin{equation}\label{eq:B-trunc-L1}
 \Ee(B_n-B_n^\circ)=O_\kappa(\log n\log\log n).
\end{equation}
In particular,
\[
 B_n-B_n^\circ=o_{L^1}((\log n)^{3/2}).
\]
\end{lemma}

\begin{proof}
By~\eqref{eq:log-cycle-tail}, the contribution from $r\leq b_n$ is
$O((\log b_n)^2)=O((\log\log n)^2)$.  By
\eqref{eq:large-log-tail}, the contribution from $r>U_n$ is
\[
 O\bigl(\log n\{1+\log(n/U_n)\}\bigr)
 =O(\log n\log\log n).
\]
This proves~\eqref{eq:B-trunc-L1}.
\end{proof}

Let
\[
 Z_n(s)=\Ee e^{sB_n^\circ},
 \qquad
 a_{n,r}=\frac{n}{r(n-r)},
 \qquad
 w_r=\log r.
\]

\begin{lemma}[Weighted switching differential equation]
\label{lem:weighted-ode}
For $s$ on the imaginary axis,
\begin{equation}\label{eq:weighted-ode}
 Z_n'(s)=H_n(s)Z_n(s)+E_n(s),
 \qquad
 H_n(s)=\sum_{b_n<r\leq U_n}a_{n,r}w_re^{sw_r},
\end{equation}
where
\begin{equation}\label{eq:E-small}
 \sup_{s\in i\mathbb R}|E_n(s)|=o((\log n)^{3/2}).
\end{equation}
\end{lemma}

\begin{proof}
Apply the marked identity~\eqref{eq:marked-one-switch} with
$\Phi(\sigma)=e^{sB_n^\circ(\sigma)}$.  If the source cycle has length
$L>r+U_n$, then $L>U_n$ and $L-r>U_n$, whereas
$b_n<r\leq U_n$.  Hence splitting the source cycle changes
$B_n^\circ$ by exactly $w_r$.  On the exceptional event
$r<L\leq r+U_n$, both exponentials have modulus one for imaginary $s$.
It follows that
\begin{align}
 &\left|\Ee[C_{n,r}e^{sB_n^\circ}]
 -a_{n,r}e^{sw_r}Z_n(s)\right|
 \notag\\
 &\quad\leq a_{n,r}\Delta_{n,r}
 +2a_{n,r}\frac1n\Ee\sum_{i:L_n(i)>r}
 R_{i,\sigma_n^r(i)}\one_{\{L_n(i)\leq r+U_n\}}.
 \label{eq:one-r-ode-error}
\end{align}
By~\eqref{eq:cutoff-range}, $r+U_n\leq2U_n\leq n/2$.  H\"older's
inequality,~\eqref{eq:selected-p}, and~\eqref{eq:root-crude} therefore give
\[
 \frac1n\Ee\sum_{i:L_n(i)>r}
 R_{i,\sigma_n^r(i)}\one_{\{L_n(i)\leq r+U_n\}}
 \leq C_\kappa\left(\frac{U_n}{n}\right)^\beta.
\]
Multiply~\eqref{eq:one-r-ode-error} by $w_r$ and sum.  On a dyadic block
$a_j\leq r<2a_j$, equation~\eqref{eq:cutoff-range} gives
$a_{n,r}\leq2/r$ for all large $n$, so
\[
 \sum_{a_j\leq r<2a_j}w_ra_{n,r}\Delta_{n,r}
 \leq C\log(2a_j)\sup_{a_j\leq r<2a_j}\Delta_{n,r}.
\]
Corollary~\ref{cor:weighted-sum} handles the sum over blocks.  Finally,
by~\eqref{eq:A0-choice},
\[
 \left(\frac{U_n}{n}\right)^\beta
 \sum_{r\leq U_n}\frac{\log r}{r}
 =O((\log n)^{2-A_0\beta})
 =o((\log n)^{3/2}).
\]
This proves~\eqref{eq:E-small}.
\end{proof}

\begin{lemma}[Deterministic exponent sums]\label{lem:det-sums}
With
\[
 a_{n,r}=\frac{n}{r(n-r)},\qquad b_n<r\leq U_n,
\]
one has
\begin{align}
 \mu_n:=\sum_{b_n<r\leq U_n}a_{n,r}\log r
 &=\frac12(\log n)^2+O(\log n\log\log n),
 \label{eq:mu-sum}\\
 v_n:=\sum_{b_n<r\leq U_n}a_{n,r}(\log r)^2
 &=\frac13(\log n)^3+O((\log n)^2\log\log n),
 \label{eq:v-sum}\\
 \sum_{b_n<r\leq U_n}a_{n,r}(\log r)^3
 &=O((\log n)^4).
 \label{eq:third-sum}
\end{align}
\end{lemma}

\begin{proof}
Use the exact decomposition
\begin{equation}\label{eq:partial-fraction}
 a_{n,r}=\frac{n}{r(n-r)}=\frac1r+\frac1{n-r}.
\end{equation}
The $1/r$ terms are estimated by the corresponding logarithmic integrals.
The lower cutoff $b_n=(\log n)^3$ and upper cutoff
$U_n=n/(\log n)^{A_0}+O(1)$ produce the displayed logarithmic errors.
For $k=1,2,3$, the $1/(n-r)$ contribution is
\[
 O\left(\frac{U_n}{n}(\log n)^k\right),
\]
which is smaller than the stated remainder.  This proves
\eqref{eq:mu-sum}--\eqref{eq:third-sum}.
\end{proof}

\begin{theorem}[Additive Erd\H{o}s--Tur\'an law]\label{thm:B-clt}
Let
\[
 \varsigma_n^2=\frac13(\log n)^3.
\]
Then
\begin{equation}\label{eq:B-clt}
 \frac{B_n-\frac12(\log n)^2}{\varsigma_n}
 \Rightarrow N(0,1).
\end{equation}
\end{theorem}

\begin{proof}
Put
\[
 \Psi_n(s)=\sum_{b_n<r\leq U_n}a_{n,r}(e^{sw_r}-1),
\]
so that $\Psi_n'=H_n$.  Variation of constants in
\eqref{eq:weighted-ode} gives
\begin{equation}\label{eq:variation}
 Z_n(s)=e^{\Psi_n(s)}
 \left[1+\int_0^s e^{-\Psi_n(z)}E_n(z)\dd z\right].
\end{equation}
For $s=it/\varsigma_n$ and $z=iu$ on the integration segment,
\begin{align*}
 -\Re\Psi_n(iu)
 &=\sum_{b_n<r\leq U_n}a_{n,r}(1-\cos(uw_r))\\
 &\leq\frac{u^2}{2}
 \sum_{b_n<r\leq U_n}a_{n,r}w_r^2.
\end{align*}
Lemma~\ref{lem:det-sums} makes the last sum $O((\log n)^3)$, so
$|e^{-\Psi_n(z)}|\leq e^{C_t}$ uniformly on the segment.  Its length is
$O(\varsigma_n^{-1})$, and~\eqref{eq:E-small} therefore makes the integral
in~\eqref{eq:variation} equal to $o(1)$.

Taylor expansion, together with Lemma~\ref{lem:det-sums}, gives for fixed
real $t$
\begin{align*}
 \Psi_n(it/\varsigma_n)
 -\frac{it}{\varsigma_n}\frac12(\log n)^2
 ={}&\frac{it}{\varsigma_n}
 \left(\mu_n-\frac12(\log n)^2\right)
 -\frac{t^2v_n}{2\varsigma_n^2}+o(1)\\
 &\longrightarrow-\frac12t^2.
\end{align*}
The cubic remainder is bounded by
\[
 \frac{C_t}{\varsigma_n^3}
 \sum_{b_n<r\leq U_n}a_{n,r}(\log r)^3
 =O((\log n)^{-1/2}).
\]
Equation~\eqref{eq:variation} and L\'evy's theorem prove
\eqref{eq:B-clt} with $B_n^\circ$ in place of $B_n$.
Lemma~\ref{lem:B-trunc} and Slutsky's theorem restore $B_n$.
\end{proof}

\section{A rate-oriented three-clock switch}

Let
\[
 p(x)=\prod_{a=1}^n\theta_a e^{-\theta_ax_a}
 \one_{\{x_a>0\ \forall a\}}
\]
be the joint clock density.  For a coordinate permutation $\tau$, define
$(T_\tau x)_a=x_{\tau(a)}$ and
\begin{equation}\label{eq:Rtau}
 R_\tau(x)=\frac{p(T_\tau x)}{p(x)}
 =\exp\left\{\sum_{a=1}^n\theta_a(x_a-x_{\tau(a)})\right\}.
\end{equation}
Then
\begin{equation}\label{eq:tau-rank}
 \sigma_{T_\tau X}=\sigma_X\circ\tau,
\end{equation}
and, by the substitution $y=T_\tau x$,
\begin{equation}\label{eq:tau-cov}
 \Ee[R_\tau(X)G(T_\tau X)]=\Ee G(X)
\end{equation}
for every nonnegative measurable $G$.

Fix three distinct labels.  Order them deterministically by decreasing rate,
with label order used to break ties, and write the resulting order as
$h_0,h_1,h_2$.  Define
\[
 \tau^*=(h_0\,h_1\,h_2).
\]
This is exactly one of the two cyclic orientations of the three labels.
Moreover,
\begin{align}
 \log R_{\tau^*}
 ={}&(\theta_{h_0}-\theta_{h_2})X_{h_0}
 -(\theta_{h_0}-\theta_{h_1})X_{h_1}
 \notag\\
 &-(\theta_{h_1}-\theta_{h_2})X_{h_2}.
 \label{eq:one-positive}
\end{align}
Thus only the maximum-rate clock has a positive coefficient.

For $r,s\geq1$ with $r+s<n$ and a source $i$ satisfying
$L_n(i)>r+s$, define
\begin{align*}
 j_i^{(1)}&=\sigma_n^r(i),&
 k_i^{(1)}&=\sigma_n^{r+s}(i),&
 \tau_i^{(1)}&=(i\,k_i^{(1)}\,j_i^{(1)}),\\
 k_i^{(2)}&=\sigma_n^s(i),&
 j_i^{(2)}&=\sigma_n^{r+s}(i),&
 \tau_i^{(2)}&=(i\,j_i^{(2)}\,k_i^{(2)}).
\end{align*}
Set
\begin{align}
 \mathcal I_{n;r,s}
 =\sum_{i:L_n(i)>r+s}\bigg[&
 \one_{\{\tau^*(i,j_i^{(1)},k_i^{(1)})=\tau_i^{(1)}\}}
 R_{\tau_i^{(1)}}(X)
 \notag\\
 &+\one_{\{\tau^*(i,j_i^{(2)},k_i^{(2)})=\tau_i^{(2)}\}}
 R_{\tau_i^{(2)}}(X)
 \bigg].
 \label{eq:I-explicit}
\end{align}
Here $\tau^*(i,j,k)$ denotes the rate-oriented cycle on the set
$\{i,j,k\}$.  When $r=s$, the two displayed terms may coincide as random
variables; they are retained with multiplicity because the target marks
order the two $r$-cycles.

\begin{proposition}[Exact three-clock factorial identity]
\label{prop:three-switch}
For $r,s\geq1$ with $r+s<n$,
\begin{equation}\label{eq:three-switch}
 rs(n-r-s)\Ee\bigl[C_{n,r}(C_{n,s}-\one_{\{r=s\}})\bigr]
 =\Ee\mathcal I_{n;r,s}.
\end{equation}
Pointwise, $\mathcal I_{n;r,s}$ is a sum of at most $2n$ nonnegative
three-clock factors, counted with multiplicity.
\end{proposition}

\begin{proof}
For a permutation $\pi$, let $A_{r,s}(\pi;i,j,k)$ be the event that $j$
lies in an $r$-cycle of $\pi$, $k$ lies in a distinct $s$-cycle, and $i$
lies in neither cycle.  The exact number of target marks is
\begin{equation}\label{eq:target-count}
 \sum_{i,j,k\ \mathrm{distinct}}
 \one_{A_{r,s}(\pi;i,j,k)}
 =rs(n-r-s)C_r(\pi)(C_s(\pi)-\one_{\{r=s\}}).
\end{equation}
For every deterministic triple $(i,j,k)$, equation~\eqref{eq:tau-cov}
gives
\begin{equation}\label{eq:triple-cov}
 \Ee\one_{A_{r,s}(\sigma_X;i,j,k)}
 =\Ee\left[
 R_{\tau^*(i,j,k)}(X)
 \one_{A_{r,s}(\sigma_X\circ\tau^*(i,j,k);i,j,k)}
 \right].
\end{equation}

If $\tau^*(i,j,k)=(i\,k\,j)$, then
\begin{equation}\label{eq:preimage-one}
 A_{r,s}(\sigma\circ(i\,k\,j);i,j,k)
 \quad\Longleftrightarrow\quad
 L_\sigma(i)>r+s,\quad
 j=\sigma^r(i),\quad
 k=\sigma^{r+s}(i).
\end{equation}
Indeed, the source cycle is traversed from $i$ to $j$ in $r$ steps, from
$j$ to $k$ in $s$ steps, and then back to $i$.  Right-composition by
$(i\,k\,j)$ replaces the three outgoing edges so as to cut off an
$r$-cycle containing $j$, an $s$-cycle containing $k$, and a residual cycle
containing $i$.  The operation is reversible.

Similarly, if $\tau^*(i,j,k)=(i\,j\,k)$, then
\begin{equation}\label{eq:preimage-two}
 A_{r,s}(\sigma\circ(i\,j\,k);i,j,k)
 \quad\Longleftrightarrow\quad
 L_\sigma(i)>r+s,\quad
 k=\sigma^s(i),\quad
 j=\sigma^{r+s}(i).
\end{equation}
Here the source cycle is traversed from $i$ to $k$ in $s$ steps and from
$k$ to $j$ in $r$ steps.  Again the right-composition cuts the source cycle
into the required three marked cycles and is reversible.

Sum~\eqref{eq:triple-cov} over all distinct triples, use
\eqref{eq:target-count}, and reorganize the right side by the source label
$i$ using~\eqref{eq:preimage-one} and~\eqref{eq:preimage-two}.  The result
is exactly~\eqref{eq:I-explicit}, proving~\eqref{eq:three-switch}.  There
are at most two terms for each source $i$.
\end{proof}

\begin{lemma}[Selected three-clock tail bound]\label{lem:three-tail}
For $\kappa>1$ and $u\geq0$,
\begin{equation}\label{eq:three-log-tail}
 \frac1n\Ee\sum_{\text{summands in }\mathcal I_{n;r,s}}
 \one_{\{\log R_{\tau^*}>u\}}
 \leq6e^{-\lambda_\kappa u},
\end{equation}
uniformly in $n,r,s$.  Consequently,
\begin{equation}\label{eq:three-first}
 \Ee\mathcal I_{n;r,s}\leq C_\kappa n.
\end{equation}
For $\kappa=1$, every factor equals one and the same conclusion holds.
\end{lemma}

\begin{proof}
Consider a selected summand and let $h$ be its maximum-rate label.  If all
three rates are equal, then $R_{\tau^*}=1$, so there is nothing to prove for
$u>0$.  Otherwise, equation~\eqref{eq:one-positive} implies
\[
 \{\log R_{\tau^*}>u\}
 \subseteq
 \left\{X_h>\frac{u}{\theta_h-\theta_{\min}}\right\}
 \subseteq
 \left\{X_h>\frac{u}{\theta_h-1}\right\}.
\]
Since $X_h\sim\Exp(\theta_h)$ and
$x/(x-1)$ is decreasing on $(1,\kappa]$,
\begin{equation}\label{eq:max-clock-tail}
 \Pp\left\{X_h>\frac{u}{\theta_h-1}\right\}
 =\exp\left\{-\frac{\theta_h}{\theta_h-1}u\right\}
 \leq e^{-\lambda_\kappa u}.
\end{equation}
Partition the selected summands according to the two orientation types in
\eqref{eq:I-explicit} and according to whether $h$ is the source, the first
endpoint, or the second endpoint.  There are at most six classes.  Within
each class, the map from the selected source $i$ to $h$ is the restriction
of one of
\[
 i\mapsto i,\qquad i\mapsto\sigma_n^r(i),\qquad
 i\mapsto\sigma_n^s(i),\qquad i\mapsto\sigma_n^{r+s}(i).
\]
Within a fixed class, the position of the maximum-rate label relative to
the source is fixed, so the map has the form $i\mapsto\sigma_n^t(i)$ for
one fixed $t\in\{0,r,s,r+s\}$.  Since $\sigma_n^t$ is a permutation of
$[n]$, its restriction to the selected source set is injective.  Reindexing
and~\eqref{eq:max-clock-tail} therefore give at most
$ne^{-\lambda_\kappa u}$ per class, proving~\eqref{eq:three-log-tail}.

Let $N_u$ be the number of selected summands with
$\log R_{\tau^*}>u$.  Since there are at most $2n$ summands,
\[
 \mathcal I_{n;r,s}
 \leq2n+\int_0^\infty e^uN_u\dd u.
\]
Taking expectations and applying~\eqref{eq:three-log-tail},
\[
 \Ee\mathcal I_{n;r,s}
 \leq2n+6n\int_0^\infty e^{-(\lambda_\kappa-1)u}\dd u
 \leq C_\kappa n.
\]
If $\kappa=1$, every summand is one and
$\mathcal I_{n;r,s}\leq2n$ directly.
\end{proof}

\begin{corollary}[Uniform two-cycle factorial bound]\label{cor:pair}
If $r+s\leq n/2$, then
\begin{equation}\label{eq:pair}
 \Ee\bigl[C_{n,r}(C_{n,s}-\one_{\{r=s\}})\bigr]
 \leq\frac{C_\kappa}{rs}.
\end{equation}
\end{corollary}

\begin{proof}
Combine~\eqref{eq:three-switch} and~\eqref{eq:three-first}, using
$n-r-s\geq n/2$.
\end{proof}

\section{The product--least-common-multiple comparison}

Let
\[
 \mathcal P_n=\{r:b_n<r\leq U_n,\ C_{n,r}>0\},
 \qquad
 O_n^\circ=\lcm\mathcal P_n,
\]
with the convention $\lcm\varnothing=1$.  Let $\Lambda$ be the von
Mangoldt function,
\[
 \Lambda(m)=
 \begin{cases}
  \log p,&m=p^a\text{ for a prime }p\text{ and }a\geq1,\\
  0,&\text{otherwise},
 \end{cases}
\]
and define, for $m\geq1$,
\[
 D_{n,m}=\sum_{\substack{b_n<r\leq U_n\\m\mid r}}C_{n,r}.
\]
The identities
\[
 \log r=\sum_{m\mid r}\Lambda(m)
\]
and
\[
 \log O_n^\circ
 =\sum_{m\leq U_n}\Lambda(m)\one_{\{D_{n,m}\geq1\}}
\]
give the exact decomposition
\begin{equation}\label{eq:prod-lcm}
 B_n^\circ-\log O_n^\circ
 =\sum_{m\leq U_n}\Lambda(m)(D_{n,m}-1)_+.
\end{equation}

\begin{proposition}[Arithmetic overlap bound]\label{prop:arith}
One has
\begin{equation}\label{eq:arith-bound}
 \Ee(B_n^\circ-\log O_n^\circ)
 =O_\kappa(\log n\log\log n).
\end{equation}
\end{proposition}

\begin{proof}
Put $y_n=(\log n)^2$.  For $m\leq y_n$, use
$(D_{n,m}-1)_+\leq D_{n,m}$.  Since $U_n\leq n/2$ for all large $n$,
Lemma~\ref{lem:crude} gives
\[
 \Ee D_{n,m}
 \leq\frac{C_\kappa}{m}\sum_{j\leq U_n/m}\frac1j
 \leq C_\kappa\frac{\log n}{m}.
\]
Chebyshev's elementary estimate~\cite{HardyWright2008}
\[
 \psi(x):=\sum_{m\leq x}\Lambda(m)=O(x)
\]
and partial summation imply
\begin{equation}\label{eq:cheb-harmonic}
 \sum_{m\leq y}\frac{\Lambda(m)}m=O(\log y).
\end{equation}
Thus the contribution of $m\leq y_n$ is
$O(\log n\log\log n)$.

For a nonnegative integer $D$,
$(D-1)_+\leq(D)_2/2$.  Moreover,
\begin{equation}\label{eq:D-factorial-expand}
 (D_{n,m})_2
 =\sum_{\substack{b_n<r,s\leq U_n\\m\mid r,\ m\mid s}}
 C_{n,r}(C_{n,s}-\one_{\{r=s\}}).
\end{equation}
By~\eqref{eq:cutoff-range}, every contributing pair satisfies
$r+s\leq2U_n\leq n/2$.  Corollary~\ref{cor:pair} therefore gives
\begin{align}
 \Ee(D_{n,m})_2
 &\leq C_\kappa
 \left(\sum_{j\leq U_n/m}\frac1{mj}\right)^2
 \notag\\
 &\leq C_\kappa\frac{\log^2(1+U_n/m)}{m^2}.
 \label{eq:D-factorial-bound}
\end{align}
A second partial summation using $\psi(x)=O(x)$ gives
\begin{equation}\label{eq:cheb-tail}
 \sum_{m>y}\frac{\Lambda(m)}{m^2}=O(1/y).
\end{equation}
Consequently, the contribution of $m>y_n$ is at most
\[
 C_\kappa(\log n)^2
 \sum_{m>y_n}\frac{\Lambda(m)}{m^2}
 =O_\kappa(1).
\]
Combining the two ranges proves~\eqref{eq:arith-bound}.
\end{proof}

Let
\[
 O_n=\ord(\sigma_n).
\]
\begin{corollary}[Product--LCM discrepancy]\label{cor:product-lcm}
One has
\begin{equation}\label{eq:full-product-lcm-L1}
 \Ee[B_n-\log O_n]=O_\kappa(\log n\log\log n).
\end{equation}
In particular,
\begin{equation}\label{eq:full-product-lcm-op}
 B_n-\log O_n=o_{L^1}((\log n)^{3/2}).
\end{equation}
\end{corollary}

\begin{proof}
The least common multiple $O_n^\circ$ divides $O_n$, while
$\log O_n\leq B_n$.  Hence the deterministic inequality
\begin{equation}\label{eq:restore-direct}
 0\leq B_n-\log O_n
 \leq(B_n-B_n^\circ)+(B_n^\circ-\log O_n^\circ)
\end{equation}
holds.  Apply Lemma~\ref{lem:B-trunc} and
Proposition~\ref{prop:arith}.
\end{proof}

\section{Proof of the main theorem}

\begin{proof}[Proof of Theorem~\ref{thm:main}]
After a common rescaling, the hypothesis of Theorem~\ref{thm:main} is
exactly~\eqref{eq:comparable}.  By Corollary~\ref{cor:product-lcm},
\[
 \frac{\log\ord(\sigma_n)-B_n}{\sqrt{\frac13(\log n)^3}}
 \longrightarrow0
\]
in $L^1$, hence in probability.  Theorem~\ref{thm:B-clt} and Slutsky's
theorem complete the proof.
\end{proof}

\begin{remark}[Quantitative refinement]
The proof deliberately retains the explicit endpoint error, the growing
truncation dependence, and the $L^1$ product--LCM discrepancy.  These
estimates form the input to a separate Berry--Esseen analysis.  No rate of
normal approximation is asserted in the present paper.
\end{remark}

\begin{thebibliography}{99}

\bibitem{ABT2003}
R. Arratia, A. D. Barbour, and S. Tavar\'e,
\newblock \emph{Logarithmic Combinatorial Structures: A Probabilistic
Approach},
\newblock EMS Monographs in Mathematics, European Mathematical Society,
Z\"urich, 2003.
\newblock \href{https://doi.org/10.4171/000}{doi:10.4171/000}.

\bibitem{BarbourTavare1994}
A. D. Barbour and S. Tavar\'e,
\newblock A rate for the Erd\H{o}s--Tur\'an law,
\newblock \emph{Combinatorics, Probability and Computing} 3 (1994), no.~2,
167--176.
\newblock \href{https://doi.org/10.1017/S0963548300001097}
{doi:10.1017/S0963548300001097}.

\bibitem{BCD2025}
J. Borga, S. Chatterjee, and P. Diaconis,
\newblock Permuton and local limits for the Luce model,
\newblock preprint, arXiv:2509.07729v2 [math.PR], 2025.
\newblock \url{https://arxiv.org/abs/2509.07729}.

\bibitem{CDK2024}
S. Chatterjee, P. Diaconis, and G. B. Kim,
\newblock Enumerative theory for the Tsetlin library,
\newblock \emph{Journal of Algebra} 655 (2024), 139--162.
\newblock \href{https://doi.org/10.1016/j.jalgebra.2023.08.009}
{doi:10.1016/j.jalgebra.2023.08.009}.

\bibitem{ErdosTuran1967}
P. Erd\H{o}s and P. Tur\'an,
\newblock On some problems of a statistical group theory. III,
\newblock \emph{Acta Mathematica Academiae Scientiarum Hungaricae} 18
(1967), 309--320.

\bibitem{GnedinIksanovMarynych2012}
A. Gnedin, A. Iksanov, and A. Marynych,
\newblock A generalization of the Erd\H{o}s--Tur\'an law for the order of
random permutation,
\newblock \emph{Combinatorics, Probability and Computing} 21 (2012), no.~5,
715--733.
\newblock \href{https://doi.org/10.1017/S0963548312000247}
{doi:10.1017/S0963548312000247}.

\bibitem{HardyWright2008}
G. H. Hardy and E. M. Wright,
\newblock \emph{An Introduction to the Theory of Numbers}, sixth edition,
\newblock Oxford University Press, Oxford, 2008.

\bibitem{LongBerry2026}
C. D. Long,
\newblock A sharp Berry--Esseen theorem and first Edgeworth expansion for the
Erd\H{o}s--Tur\'an law of comparable-rate Luce permutations,
\newblock research supplement, 2026.

\bibitem{LongLuce2026}
C. D. Long,
\newblock Boundary, terminal, diagonal, and cycle limit laws for Luce
permutations,
\newblock manuscript under referee revision, 2026.

\bibitem{LongSwitch2026}
C. D. Long,
\newblock Mesoscopic cycles of comparable-rate Luce permutations: a
switching proof of Conjecture~13.1,
\newblock manuscript, 2026.

\bibitem{Luce1959}
R. D. Luce,
\newblock \emph{Individual Choice Behavior: A Theoretical Analysis},
\newblock John Wiley \& Sons, New York, 1959.

\bibitem{StormZeindler2015}
J. Storm and D. Zeindler,
\newblock The order of large random permutations with cycle weights,
\newblock \emph{Electronic Journal of Probability} 20 (2015), Paper No.~126,
34 pp.
\newblock \href{https://doi.org/10.1214/EJP.v20-4331}
{doi:10.1214/EJP.v20-4331}.

\end{thebibliography}

\end{document}